% Source: https://artofproblemsolving.com/community/c3h2267362_remainder_theorem
% Stable identity: aops-topic-2267362-q01
\textbf{(a)} By the remainder theorem, $V(c)$ is the remainder when $V(t)$ is divided by $t-c$. Synthetic division with the coefficients $1,-19,108,-180,0$:
\begin{itemize}
\item $c=1$: $1,\ -18,\ 90,\ -90,\ -90$, so $V(1)=-90$;
\item $c=5$: $1,\ -14,\ 38,\ 10,\ 50$, so $V(5)=50$;
\item $c=8$: $1,\ -11,\ 20,\ -20,\ -160$, so $V(8)=-160$.
\end{itemize}
The speeds are $200-90=\boxed{110}$ mph at $t=1$, $200+50=\boxed{250}$ mph at $t=5$, and $200-160=\boxed{40}$ mph at $t=8$.

\textbf{(b)} The speed is $200$ mph exactly when $V(t)=0$. Factor out $t$: $V(t)=t\,(t^3-19t^2+108t-180)$. Trying divisors of $180$, $t=3$ is a root of the cubic ($27-171+324-180=0$); synthetic division by $t-3$ gives $t^2-16t+60=(t-6)(t-10)$. So
\[V(t)=t(t-3)(t-6)(t-10).\]
The roots in $0<t<10.3$ are $t=3,6,10$, so his speed was $200$ mph \textbf{three times}, at $t=3$, $t=6$ and $t=10$ seconds.
